\documentclass[11pt, a4paper]{article}

% === LuaLaTeX 패키지 ===
\usepackage{fontspec}
\usepackage{kotex}
\usepackage{emoji}
\setemojifont{Noto Color Emoji}
\directlua{
  luaotfload.add_fallback("emojifb", {
    "NotoColorEmoji:mode=harf;"
  })
}
\setmainfont{Noto Sans CJK KR}[RawFeature={fallback=emojifb}]
\setsansfont{Noto Sans CJK KR}[RawFeature={fallback=emojifb}]
\setmonofont{Noto Sans Mono CJK KR}[Scale=0.85]

% === 레이아웃 ===
\usepackage[margin=2.3cm, top=2.5cm, bottom=2.5cm]{geometry}
\usepackage{setspace}\onehalfspacing

% === 패키지 ===
\usepackage{graphicx}
\usepackage{booktabs}
\usepackage{array}
\usepackage{longtable}
\usepackage{xcolor}
\usepackage{tcolorbox}
\usepackage{enumitem}
\usepackage{fancyhdr}
\usepackage{titlesec}
\usepackage{hyperref}
\usepackage{float}

% === 한글 캡션 ===
\renewcommand{\contentsname}{목차}
\renewcommand{\tablename}{표}
\renewcommand{\figurename}{그림}

% === 색상 ===
\definecolor{mainblue}{RGB}{21, 101, 192}
\definecolor{accentgreen}{RGB}{46, 125, 50}
\definecolor{warningorange}{RGB}{230, 81, 0}
\definecolor{lightbg}{RGB}{245, 245, 245}
\definecolor{tipbg}{RGB}{232, 245, 233}
\definecolor{corebg}{RGB}{227, 242, 253}

% === tcolorbox ===
\tcbuselibrary{skins, breakable}
\newtcolorbox{corebox}[1][]{
  colback=corebg, colframe=mainblue, boxrule=1pt, arc=4pt,
  fonttitle=\bfseries\color{mainblue}, title={#1}, breakable
}

% === 헤더/푸터 ===
\pagestyle{fancy}\fancyhf{}
\fancyhead[L]{\small\textcolor{mainblue}{사물인터넷과 빅데이터}}
\fancyhead[R]{\small\textcolor{gray}{2주차: 강의개요}}
\fancyfoot[C]{\thepage}
\renewcommand{\headrulewidth}{0.5pt}

% === 섹션 스타일 ===
\titleformat{\section}{\Large\bfseries\color{mainblue}}{\thesection.}{0.5em}{}
\titleformat{\subsection}{\large\bfseries\color{mainblue!80}}{\thesubsection}{0.5em}{}
\titleformat{\subsubsection}{\normalsize\bfseries\color{mainblue!60}}{\thesubsubsection}{0.5em}{}
\hypersetup{colorlinks=true, linkcolor=mainblue, urlcolor=mainblue!70}

% ============================================================
\begin{document}

% === 표지 ===
\begin{titlepage}
\centering
\vspace*{3cm}
{\Huge\bfseries\color{mainblue} 사물인터넷과 빅데이터\par}
\vspace{0.5cm}
{\LARGE 2주차 강의개요\par}
\vspace{1.5cm}
{\Huge\bfseries 바이브코딩 소개\par}
\vspace{2cm}
{\large 청주교육대학교 컴퓨터교육과\par}
\vspace{0.3cm}
{\large 2026학년도 1학기\par}
\vspace{3cm}
{\normalsize 문서 버전: v1.0.0.0\par}
{\normalsize 작성일: 2026-02-19\par}
\end{titlepage}

% ============================================================
\section{기본 정보}
% ============================================================

\begin{table}[H]\centering
\caption{2주차 기본 정보}
\begin{tabular}{ll}\toprule
\textbf{항목} & \textbf{내용} \\\midrule
주차 & 2주차 \\
주제 & 바이브코딩 소개 \\
학습 목표 & 바이브코딩의 개념을 이해하고, 다양한 센서를 탐색한다. \\
운영 방식 & 대면 수업 (100분) \\
바이브코딩 & 활용 \\
선수 학습 & 1주차: 마이크로비트 LED 이름 표시 체험 \\
\bottomrule
\end{tabular}
\end{table}

% ============================================================
\section{수업 흐름}
% ============================================================

\begin{table}[H]\centering
\caption{수업 시간 배분}
\begin{tabular}{ccll}\toprule
\textbf{순서} & \textbf{시간} & \textbf{활동} & \textbf{유형} \\\midrule
1 & 20분 & 바이브코딩 개념 설명 & 강의 \\
2 & 20분 & Claude 사용법 시연 & 시연 \\
3 & 20분 & 바이브코딩 체험: 자기소개 홈페이지 만들기 & 실습 \\
4 & 30분 & 센서 탐색: 내장 센서 + Grove 확장 센서 & 강의+탐색 \\
5 & 10분 & 프로젝트 주제 브레인스토밍 안내 & 안내 \\
\bottomrule
\end{tabular}
\end{table}

% ============================================================
\section{학습 내용 개요}
% ============================================================

\subsection{바이브코딩이란? (활동 1)}

\begin{itemize}[leftmargin=1.5em]
\item 바이브코딩의 정의: ``코드를 짜지 않고, AI에게 말로 설명하여 코드를 만드는 방법''
\item 전통 코딩 vs 바이브코딩 비교
\item 바이브코딩의 장점과 한계
\item 이 수업에서 바이브코딩을 사용하는 이유: 코드 문법이 아닌 \textbf{문제 해결 과정}에 집중
\end{itemize}

\subsection{Claude 사용법 (활동 2)}

\begin{itemize}[leftmargin=1.5em]
\item Claude 접속 및 기본 인터페이스 안내
\item 좋은 프롬프트 작성법 (구체적으로, 단계적으로, 예시 포함)
\item 프롬프트 → 결과물 → 수정 요청의 반복 과정
\item Artifact(코드 결과물) 확인 및 다운로드 방법
\item 대화목록 PDF 생성 방법 (과제 제출용)
\end{itemize}

\subsection{바이브코딩 체험 (활동 3)}

\begin{itemize}[leftmargin=1.5em]
\item 실습 과제: Claude에게 자기소개 홈페이지 만들어달라고 요청
\item 첫 프롬프트 작성 → 결과 확인 → 수정 요청 → 완성
\item HTML 파일 다운로드 및 브라우저에서 확인
\item 체험 후 소감: ``코딩을 몰라도 웹페이지를 만들 수 있다!''
\end{itemize}

\subsection{센서 탐색 (활동 4)}

\subsubsection{마이크로비트 V2 내장 센서 (회로 연결 없이 사용)}

\begin{itemize}[leftmargin=1.5em]
\item 온도 센서, 마이크(소음), 조도 센서, 가속도계, 나침반, 터치 로고
\item 각 센서의 측정 범위, 특징, 한계
\item 1주차 복습 연결: LED 이름 표시에서 → 센서 데이터 활용으로 확장
\end{itemize}

\subsubsection{Grove 확장 센서 (I2C 플러그앤플레이)}

\begin{itemize}[leftmargin=1.5em]
\item Grove Shield 연결 방식 소개 (케이블 꽂기만 하면 됨)
\item I2C 통신 개념 (간단히: ``주소로 센서를 구분하는 버스'')
\item 주요 Grove I2C 센서: 온습도(SHT40), 조도(TSL2561), 색상(TCS34725), 거리(VL53L0X), 기압(BMP280), 환경(BME680), 가속도(LIS3DHTR)
\item 센서 조합의 다양성: 내장 + 외장 조합으로 848가지 가능
\end{itemize}

\subsubsection{센서 시뮬레이터 시연 (참고 도구)}

\begin{itemize}[leftmargin=1.5em]
\item 센서 시뮬레이터 v9: 센서 슬롯 배치 및 가상 데이터 확인
\item 센서 검색 도구: ``교실 온도와 습도를 측정하고 싶어요'' → 센서 추천
\item I2C 연결도: 센서별 주소와 연결 구조 시각화
\end{itemize}

\subsection{프로젝트 브레인스토밍 안내 (활동 5)}

\begin{itemize}[leftmargin=1.5em]
\item 다음 주(3주차) 프로젝트 기획 예고
\item ``학교 현장에서 무엇을 측정할 수 있을까?'' 생각해오기
\item 프로젝트 예시 간단 소개: 교실 환경, 복도 소음, 급식실 대기, 운동장 활동량
\end{itemize}

% ============================================================
\section{교수·학생 활동}
% ============================================================

\subsection{교수 활동}

\begin{itemize}[leftmargin=1.5em]
\item 바이브코딩 개념 설명 및 시연
\item Claude 사용법 시연 (프롬프트 작성 → 결과 → 수정 과정)
\item 센서 종류 소개 (내장 + Grove)
\item 센서 시뮬레이터 시연
\end{itemize}

\subsection{학생 활동}

\begin{itemize}[leftmargin=1.5em]
\item 바이브코딩 체험 (자기소개 홈페이지 만들기)
\item Claude와 대화하며 결과물 개선
\item 센서 탐색 (시뮬레이터로 센서 조합 확인)
\item 프로젝트 아이디어 구상 시작
\end{itemize}

% ============================================================
\section{과제}
% ============================================================

\begin{table}[H]\centering
\caption{과제 안내}
\begin{tabular}{ll}\toprule
\textbf{항목} & \textbf{내용} \\\midrule
과제명 & 바이브코딩 체험 + 센서 탐색 \\
배점 & 5\% \\
제출물 & 결과물(HTML) + 대화목록 전문(PDF) + 대화목록 요약본(PDF) \\
제출 기한 & 일요일 자정 (LMS) \\
대화목록 생성 & Claude에게 프롬프트로 요청하여 자동 생성 \\
\bottomrule
\end{tabular}
\end{table}

% ============================================================
\section{수업 도구 및 자료}
% ============================================================

\subsection{강의 자료 (수업 중 사용)}

\begin{table}[H]\centering
\begin{tabular}{ll}\toprule
\textbf{자료} & \textbf{용도} \\\midrule
강의교재 (PDF) & 수업 내용 교재 \\
강의슬라이드 (PDF) & 발표 자료 \\
\bottomrule
\end{tabular}
\end{table}

\subsection{참고 자료 (시연·탐색용)}

\begin{table}[H]\centering
\begin{tabular}{lll}\toprule
\textbf{도구} & \textbf{버전} & \textbf{용도} \\\midrule
센서 시뮬레이터 & v9.1.0.0 & 센서 슬롯 배치, 가상 데이터 확인 \\
센서 검색 도구 & v1.0.2.0 & 자연어로 센서 추천 \\
I2C 연결도 & v1.0.2.0 & 센서 주소·연결 구조 시각화 \\
데이터로거 & v1.0.2.0 & 데이터 수집 구조 소개 \\
집중도 분석 프로그램 & v2.2.0.0 & 교수자 예시 프로젝트 시연 \\
\bottomrule
\end{tabular}
\end{table}

\subsection{외부 도구}

\begin{table}[H]\centering
\begin{tabular}{lll}\toprule
\textbf{도구} & \textbf{URL} & \textbf{용도} \\\midrule
Claude & claude.ai & 바이브코딩 실습 \\
MakeCode & makecode.microbit.org & 참고 (1주차 복습) \\
\bottomrule
\end{tabular}
\end{table}

% ============================================================
\section{그림 목록}
% ============================================================

강의교재 및 슬라이드에 포함할 그림 목록입니다.

\begin{table}[H]\centering
\caption{그림 목록}
\small
\begin{tabular}{clll}\toprule
\textbf{\#} & \textbf{그림} & \textbf{용도} & \textbf{출처} \\\midrule
1 & 전통 코딩 vs 바이브코딩 비교 & 바이브코딩 개념 설명 & 신규 제작 (SVG) \\
2 & Claude 인터페이스 화면 구성 & Claude 사용법 안내 & 신규 제작 (SVG) \\
3 & 프롬프트 → 결과 → 수정 흐름도 & 바이브코딩 작업 흐름 & 신규 제작 (SVG) \\
4 & 마이크로비트 V2 내장 센서 배치도 & 내장 센서 소개 & 신규 제작 (SVG) \\
5 & Grove Shield 연결 다이어그램 & Grove 확장 방식 소개 & 신규 제작 (SVG) \\
6 & I2C 버스 개념도 & I2C 통신 원리 설명 & 신규 제작 (SVG) \\
7 & Grove I2C 센서 카탈로그 & 센서 옵션 소개 & 신규 제작 (SVG) \\
8 & 센서 시뮬레이터 화면 캡처 & 시뮬레이터 소개 & 기존 스크린샷 \\
9 & 센서 검색 도구 화면 캡처 & 검색 도구 소개 & 기존 스크린샷 \\
10 & I2C 연결도 화면 캡처 & 연결도 소개 & 기존 스크린샷 \\
11 & 프로젝트 예시 개요도 & 브레인스토밍 안내 & 신규 제작 (SVG) \\
\bottomrule
\end{tabular}
\end{table}

% ============================================================
\section{핵심 메시지 연결}
% ============================================================

\begin{table}[H]\centering
\caption{핵심 메시지와 2주차 적용}
\begin{tabular}{p{5cm}p{7.5cm}}\toprule
\textbf{강좌 핵심 메시지} & \textbf{2주차 적용} \\\midrule
``데이터는 원인을 말하지 않는다'' & 센서 탐색 시 ``센서가 측정하는 것 vs 우리가 알고 싶은 것''의 차이 언급 \\
``AI는 판단하지 않는다, 계산할 뿐이다'' & Claude가 코드를 `이해'하는 게 아니라 패턴으로 생성한다는 점 \\
``코딩을 의심하는 방법을 가르친다'' & 바이브코딩 결과물을 무조건 신뢰하지 말고 확인하는 습관 \\
\bottomrule
\end{tabular}
\end{table}

% ============================================================
\section*{참고자료}
\addcontentsline{toc}{section}{참고자료}

\begin{table}[H]\centering
\begin{tabular}{clll}\toprule
\textbf{\#} & \textbf{자료명} & \textbf{유형} & \textbf{비고} \\\midrule
1 & 강의계획서 v1.5.3.0 & 문서 & 2주차 상세 계획 \\
2 & 센서 시뮬레이터 v9.1.0.0 & HTML & 참고 도구 \\
3 & 센서 검색 도구 v1.0.2.0 & HTML & 참고 도구 \\
4 & I2C 연결도 v1.0.2.0 & HTML & 참고 도구 \\
\bottomrule
\end{tabular}
\end{table}

% ============================================================
\section*{생성/수정 이력}
\addcontentsline{toc}{section}{생성/수정 이력}

\begin{table}[H]\centering
\begin{tabular}{llllll}\toprule
\textbf{버전} & \textbf{날짜} & \textbf{시간} & \textbf{수준} & \textbf{변경 내용} & \textbf{작성자} \\\midrule
v1.0.0.0 & 2026-02-19 & 21:00 & 최초 & 강의계획서 기반 2주차 강의개요 초안 & Claude \\
\bottomrule
\end{tabular}
\end{table}

\end{document}
